\section{WL refinement and homomorphism tests}\label{sec:wl-tests}

We use the refinement conventions and parameter \(t\) from
Section~\ref{sec:introduction}; throughout \(k\ge4\).
We will use homomorphism counts in two ways to connect the construction to WL.
Equality of all counts from sources with small bags will prove
equivalence. A difference for one explicit template will prove
inequivalence, already after two joint rounds or one separate round.

A graph has \emph{bag size at most \(t\)} if it has a tree decomposition
on a finite tree with bags of size at most \(t\): the bags containing any vertex form a
nonempty connected subtree, and some bag contains both ends of each
edge. Equivalently, its treewidth is at most \(t-1\).
Write \(\hom(F,X)\) for the number of homomorphisms from \(F\) to \(X\),
that is, vertex maps preserving adjacency.

\begin{definition}\label{def:template}
The template \(T=T_t\) has main vertices \([t]\) forming a clique.
For each \(i<j\) and \(\ell\in[t]\setminus\{i,j\}\), it has a helper
\(h_{ij;\ell}\) adjacent precisely to \(i,j,\ell\).
\end{definition}
Its main clique bag, with one attached bag
\(\{i,j,\ell,h_{ij;\ell}\}\) per helper, gives bag size at most \(t\),
since \(t\ge4\).
Every main triple supports three helpers. Consequently
\begin{equation}\label{eq:template-size}
 |V(T)|=t+3\binom t3,\qquad
 \deg_T(i)=t-1+3\binom{t-1}{2},\qquad
 \deg_T(h_{ij;\ell})=3.
\end{equation}
In particular \(T\) is connected and every vertex has positive degree.

We first establish the preservation statement for arbitrary sources
with bags of size at most \(t\), then prove the detection statement
for \(T_t\).
The homomorphism characterization of WL is established in
\citet{Dvorak2010,DGR2018}. We give a self-contained adaptation of the
rooted-count argument for the preservation direction.

\begin{lemma}\label{lem:hom-implies-wl}
If \(\hom(F,X_0)=\hom(F,X_1)\) for every graph \(F\) of bag size at most
\(t\), then \(X_0,X_1\) are \(k\)-WL equivalent in the corresponding
convention. This also holds for the disjoint-union implementation.
\end{lemma}

\begin{proof}
We realize WL updates by operations on rooted counts, then recover
color indicators on the fixed finite targets by interpolation.
Equip \(F\) with a list \(\boldsymbol\ell=(\ell_1,\ldots,\ell_k)\) of
roots, allowing repetitions, all contained in a designated bag.
Let \(h_F(X,\mathbf x)\) count homomorphisms sending \(\ell_i\) to \(x_i\),
and let \(\mathcal A_t\) be the rational span of these functions for
sources with bag size at most \(t\).
Distinct isolated roots give the constant \(1\); repeating a root tests
equality, and joining two roots tests adjacency. Complements and products
therefore give atomic-type indicators.

Here is the decomposition invariant behind the algebra.
To glue rooted graphs along an interface, attach their designated bags
to a central bag of formal interface vertices and identify each root
with its assigned interface vertex. In each operand, every merged
vertex is a root, so its occurrence subtrees meet the designated bag.
Their quotient occurrences connect through the central bag. Other
vertices retain their connected occurrence subtrees, and taking images
of bags never increases their sizes. Parallel edges may be collapsed;
a loop gives the zero function on simple targets. The central bag
contains the output roots, so this invariant persists under successive
operations. In particular, gluing corresponding roots realizes products
in \(\mathcal A_t\), with an interface of size at most \(k\).

For joint refinement, the same construction gives closure under
\[
 (f_1,\ldots,f_k)\longmapsto
 \left[(X,\mathbf x)\longmapsto
       \sum_{y\in V(X)}\prod_{i=1}^k
                      f_i(X,\mathbf x[i\leftarrow y])\right].
\]
For rooted counts, use a central bag of formal vertices
\(u_1,\ldots,u_k,y\), attach operand \(i\) at its replacement list, and
retain \(u_1,\ldots,u_k\) as output roots. This uses \(k+1=t\) vertices;
multilinearity handles their linear combinations.
Repeated roots can force \(y=u_j\) or merge old positions, precisely
enforcing the corresponding equalities in the summation.

For separate refinement only
\(f\mapsto\sum_y f(X,\mathbf x[i\leftarrow y])\) is needed.
Use adjacent bags
\(O=\{u_1,\ldots,u_k\}\) and
\(R=\{u_j:j\ne i\}\cup\{y\}\), and attach the operand's root bag to \(R\).
Each has size at most \(k=t\).
After the root identifications, a merged class meeting both \(O\)
and the operand contains a surviving \(u_j\), so it also meets \(R\).
Thus no occurrence gap is created, even when \(y=u_j\).
The old \(u_i\) remains a pinned output root in \(O\).
A homomorphism of the quotient specifies exactly one allowed image
of \(y\) and one operand homomorphism: an identification restricts that
image, rather than contributing an additional free choice.

Fix the two finite targets. At each round, every color indicator
restricted to tuples of these targets agrees with a function in
\(\mathcal A_t\). Indeed the update operations give the
multiplicity of each color vector (or each coordinate color).
If such a multiplicity \(M\) has attained values \(S\), its value-\(q\)
indicator is the polynomial
\(\prod_{s\in S\setminus\{q\}}(M-s)/(q-s)\), for \(q\in S\).
Products with the old-color indicator give every new-color indicator.
The coefficients are common to the two targets. Finally,
\(\sum_{\mathbf x\in V(X)^k}h_F(X,\mathbf x)=\hom(F,X)\):
each homomorphism specifies one root-image tuple, even with repeated
roots. The hypothesis therefore equates all color histograms.

For the union implementation, perform the indicator induction on
\emph{all} tuples of \(U=X_0\sqcup X_1\), including mixed tuples.
The interpolation coefficients are then fixed before either pure
restriction is summed. For the final quotient source \(F\), let
\(\mathcal L\) be its components containing roots and \(\mathcal R\)
its rootless components. The exact identity is
\begin{equation}\label{eq:pure-tuple-factorization}
 \sum_{\mathbf x\in V(X_j)^k}h_F(U,\mathbf x)
   =\prod_{Q\in\mathcal L}\hom(Q,X_j)
      \prod_{Q\in\mathcal R}\hom(Q,U).
\end{equation}
Connectedness forces each rooted component into \(X_j\); rootless
components are unrestricted. Every map still determines one tuple.
Different source components may share target images, so their choices
remain independent. Components inherit the bag bound, the first
product agrees for \(j=0,1\), and the second is common.
Applying this identity to each term of a color indicator proves the
claim.
\end{proof}

It remains to show that a count difference for the template \(T_t\)
is visible in the claimed number of rounds. Its helpers have only
three neighbors, so their possible images can be counted from
triple common-neighbor counts.

\begin{lemma}\label{lem:test-detected}
The round-two color histogram in joint \(k\)-WL determines
\(\hom(T_{k+1},X)\). The round-one histogram in separate \(k\)-WL
determines \(\hom(T_k,X)\). The same statements hold for pure-tuple
histograms obtained by refining a disjoint union.
\end{lemma}

\begin{proof}
Let \(A_X(v,w)=\ind{vw\in E(X)}\).
For main images \(\mathbf z=(z_1,\ldots,z_t)\), put
\(q_B(\mathbf z)=|\bigcap_{b\in B}N_X(z_b)|\), for triples \(B\subseteq[t]\).
Their number of extensions to \(T\) is
\[
 E_X(\mathbf z)=
       \prod_{a<b} A_X(z_a,z_b)
       \prod_{B\in\binom{[t]}3}q_B(\mathbf z)^3.
\]
The clique factor rejects repeated main images. Helper images may
coincide, so their choices are independent and the displayed product
is exact.

First consider joint refinement. A round-one color of a
\(k\)-tuple determines each triple common-neighbor count on its positions.
For one candidate \(w\), its atomic replacement vector reveals adjacency
to position \(p\) in any component replacing a coordinate other than \(p\).
Thus that one vector tests simultaneous adjacency to the whole triple;
summing its multiplicities gives the count.

Now fix \(\mathbf x=(z_1,\ldots,z_k)\) and set \(z_{k+1}=y\).
Its old round-one color and the single vector
\[
 \bigl(C_1(\mathbf x[1\leftarrow y]),\ldots,
       C_1(\mathbf x[k\leftarrow y])\bigr)
\]
give the round-one colors of all \(k\)-position faces of these \(k+1\)
main positions. The old tuple omits \(k+1\); replacement face \(i\)
omits \(i\), with its coordinate \(i\) representing position \(k+1\).
Every pair and triple lies in a face, since \(k\ge4\).
Reading it from, for example, the face omitting the least position
outside it determines every factor of \(E_X(\mathbf x,y)\).
All faces in this vector refer to the same \(y\).
The round-two multiset therefore determines
\(\sum_y E_X(\mathbf x,y)\); summing over tuple colors gives
\(\hom(T_{k+1},X)\).

For separate refinement, all \(t=k\) main images are already fixed.
For a triple \(B\), choose \(i\notin B\). If \(M_i(\sigma)\) is the
multiplicity of atomic type \(\sigma\) in the \(i\)-th replacement
multiset, then
\[
 q_B(\mathbf x)=\sum_\sigma M_i(\sigma)
                    \prod_{p\in B} A_\sigma(i,p),
\]
where \(A_\sigma(i,p)\) is its adjacency bit.
Thus the round-one color determines \(E_X(\mathbf x)\), and its
histogram determines \(\hom(T_k,X)\).

Finally take a pure tuple in \(X_j\) while refining \(U=X_0\sqcup X_1\).
In the joint case, a last main image outside \(X_j\) makes the
clique factor zero. Once all main images lie in \(X_j\), every common
neighbor of a triple lies there too. Hence the same formulas, evaluated
in \(U\), count precisely the desired extensions into \(X_j\).
\end{proof}
